\documentclass[a4paper,12pt,twoside]{memoir}

% Castellano
\usepackage[spanish,es-tabla]{babel}
\selectlanguage{spanish}
\usepackage[utf8]{inputenc}
\usepackage[T1]{fontenc}
\usepackage{lmodern} % Scalable font
\usepackage{microtype}
\usepackage{placeins}
\usepackage{titlesec}

\RequirePackage{booktabs}
\RequirePackage[table]{xcolor}
\RequirePackage{xtab}
\RequirePackage{multirow}

\setcounter{tocdepth}{2}
% Links
\PassOptionsToPackage{hyphens}{url}\usepackage[colorlinks]{hyperref}
\hypersetup{
	allcolors = {red}
}

% Ecuaciones
\usepackage{amsmath}

% Rutas de fichero / paquete
\newcommand{\ruta}[1]{{\sffamily #1}}

% Párrafos
\nonzeroparskip

% Huérfanas y viudas
\widowpenalty100000
\clubpenalty100000

% Imágenes

% Comando para insertar una imagen en un lugar concreto.
% Los parámetros son:
% 1 --> Ruta absoluta/relativa de la figura
% 2 --> Texto a pie de figura
% 3 --> Tamaño en tanto por uno relativo al ancho de página
\usepackage{graphicx}
\newcommand{\imagen}[3]{
	\begin{figure}[!h]
		\centering
		\includegraphics[width=#3\textwidth]{#1}
		\caption{#2}\label{fig:#1}
	\end{figure}
	\FloatBarrier
}

% Comando para insertar una imagen sin posición.
% Los parámetros son:
% 1 --> Ruta absoluta/relativa de la figura
% 2 --> Texto a pie de figura
% 3 --> Tamaño en tanto por uno relativo al ancho de página
\newcommand{\imagenflotante}[3]{
	\begin{figure}
		\centering
		\includegraphics[width=#3\textwidth]{#1}
		\caption{#2}\label{fig:#1}
	\end{figure}
}

% El comando \figura nos permite insertar figuras comodamente, y utilizando
% siempre el mismo formato. Los parametros son:
% 1 --> Porcentaje del ancho de página que ocupará la figura (de 0 a 1)
% 2 --> Fichero de la imagen
% 3 --> Texto a pie de imagen
% 4 --> Etiqueta (label) para referencias
% 5 --> Opciones que queramos pasarle al \includegraphics
% 6 --> Opciones de posicionamiento a pasarle a \begin{figure}
\newcommand{\figuraConPosicion}[6]{%
  \setlength{\anchoFloat}{#1\textwidth}%
  \addtolength{\anchoFloat}{-4\fboxsep}%
  \setlength{\anchoFigura}{\anchoFloat}%
  \begin{figure}[#6]
    \begin{center}%
      \Ovalbox{%
        \begin{minipage}{\anchoFloat}%
          \begin{center}%
            \includegraphics[width=\anchoFigura,#5]{#2}%
            \caption{#3}%
            \label{#4}%
          \end{center}%
        \end{minipage}
      }%
    \end{center}%
  \end{figure}%
}

%
% Comando para incluir imágenes en formato apaisado (sin marco).
\newcommand{\figuraApaisadaSinMarco}[5]{%
  \begin{figure}%
    \begin{center}%
    \includegraphics[angle=90,height=#1\textheight,#5]{#2}%
    \caption{#3}%
    \label{#4}%
    \end{center}%
  \end{figure}%
}
% Para las tablas
\newcommand{\otoprule}{\midrule [\heavyrulewidth]}
%
% Nuevo comando para tablas pequeñas (menos de una página).
\newcommand{\tablaSmall}[5]{%
 \begin{table}
  \begin{center}
   \rowcolors {2}{gray!35}{}
   \begin{tabular}{#2}
    \toprule
    #4
    \otoprule
    #5
    \bottomrule
   \end{tabular}
   \caption{#1}
   \label{tabla:#3}
  \end{center}
 \end{table}
}

%
% Nuevo comando para tablas pequeñas (menos de una página).
\newcommand{\tablaSmallSinColores}[5]{%
 \begin{table}[H]
  \begin{center}
   \begin{tabular}{#2}
    \toprule
    #4
    \otoprule
    #5
    \bottomrule
   \end{tabular}
   \caption{#1}
   \label{tabla:#3}
  \end{center}
 \end{table}
}

\newcommand{\tablaApaisadaSmall}[5]{%
\begin{landscape}
  \begin{table}
   \begin{center}
    \rowcolors {2}{gray!35}{}
    \begin{tabular}{#2}
     \toprule
     #4
     \otoprule
     #5
     \bottomrule
    \end{tabular}
    \caption{#1}
    \label{tabla:#3}
   \end{center}
  \end{table}
\end{landscape}
}

%
% Nuevo comando para tablas grandes con cabecera y filas alternas coloreadas en gris.
\newcommand{\tabla}[6]{%
  \begin{center}
    \tablefirsthead{
      \toprule
      #5
      \otoprule
    }
    \tablehead{
      \multicolumn{#3}{l}{\small\sl continúa desde la página anterior}\\
      \toprule
      #5
      \otoprule
    }
    \tabletail{
      \hline
      \multicolumn{#3}{r}{\small\sl continúa en la página siguiente}\\
    }
    \tablelasttail{
      \hline
    }
    \bottomcaption{#1}
    \rowcolors {2}{gray!35}{}
    \begin{xtabular}{#2}
      #6
      \bottomrule
    \end{xtabular}
    \label{tabla:#4}
  \end{center}
}

%
% Nuevo comando para tablas grandes con cabecera.
\newcommand{\tablaSinColores}[6]{%
  \begin{center}
    \tablefirsthead{
      \toprule
      #5
      \otoprule
    }
    \tablehead{
      \multicolumn{#3}{l}{\small\sl continúa desde la página anterior}\\
      \toprule
      #5
      \otoprule
    }
    \tabletail{
      \hline
      \multicolumn{#3}{r}{\small\sl continúa en la página siguiente}\\
    }
    \tablelasttail{
      \hline
    }
    \bottomcaption{#1}
    \begin{xtabular}{#2}
      #6
      \bottomrule
    \end{xtabular}
    \label{tabla:#4}
  \end{center}
}

%
% Nuevo comando para tablas grandes sin cabecera.
\newcommand{\tablaSinCabecera}[5]{%
  \begin{center}
    \tablefirsthead{
      \toprule
    }
    \tablehead{
      \multicolumn{#3}{l}{\small\sl continúa desde la página anterior}\\
      \hline
    }
    \tabletail{
      \hline
      \multicolumn{#3}{r}{\small\sl continúa en la página siguiente}\\
    }
    \tablelasttail{
      \hline
    }
    \bottomcaption{#1}
  \begin{xtabular}{#2}
    #5
   \bottomrule
  \end{xtabular}
  \label{tabla:#4}
  \end{center}
}



\definecolor{cgoLight}{HTML}{EEEEEE}
\definecolor{cgoExtralight}{HTML}{FFFFFF}

%
% Nuevo comando para tablas grandes sin cabecera.
\newcommand{\tablaSinCabeceraConBandas}[5]{%
  \begin{center}
    \tablefirsthead{
      \toprule
    }
    \tablehead{
      \multicolumn{#3}{l}{\small\sl continúa desde la página anterior}\\
      \hline
    }
    \tabletail{
      \hline
      \multicolumn{#3}{r}{\small\sl continúa en la página siguiente}\\
    }
    \tablelasttail{
      \hline
    }
    \bottomcaption{#1}
    \rowcolors[]{1}{cgoExtralight}{cgoLight}

  \begin{xtabular}{#2}
    #5
   \bottomrule
  \end{xtabular}
  \label{tabla:#4}
  \end{center}
}



\graphicspath{ {./img/} }

% Capítulos
\chapterstyle{bianchi}
\newcommand{\capitulo}[2]{
	\setcounter{chapter}{#1}
	\setcounter{section}{0}
	\setcounter{figure}{0}
	\setcounter{table}{0}
	\chapter*{#2}
	\addcontentsline{toc}{chapter}{#2}
	\markboth{#2}{#2}
}

% Apéndices
\renewcommand{\appendixname}{Apéndice}
\renewcommand*\cftappendixname{\appendixname}

\newcommand{\apendice}[1]{
	%\renewcommand{\thechapter}{A}
	\chapter{#1}
}

\renewcommand*\cftappendixname{\appendixname\ }

% Formato de portada
\makeatletter
\usepackage{xcolor}
\newcommand{\tutor}[1]{\def\@tutor{#1}}
\newcommand{\course}[1]{\def\@course{#1}}
\definecolor{cpardoBox}{HTML}{E6E6FF}
\def\maketitle{
  \null
  \thispagestyle{empty}
  % Cabecera ----------------
\noindent\includegraphics[width=\textwidth]{cabecera}\vspace{1cm}%
  \vfill
  % Título proyecto y escudo informática ----------------
  \colorbox{cpardoBox}{%
    \begin{minipage}{.8\textwidth}
      \vspace{.5cm}\Large
      \begin{center}
      \textbf{TFG del Grado en Ingeniería Informática}\vspace{.6cm}\\
      \textbf{\LARGE\@title{}}
      \end{center}
      \vspace{.2cm}
    \end{minipage}

  }%
  \hfill\begin{minipage}{.20\textwidth}
    \includegraphics[width=\textwidth]{escudoInfor}
  \end{minipage}
  \vfill
  % Datos de alumno, curso y tutores ------------------
  \begin{center}%
  {%
    \noindent\LARGE
    Presentado por \@author{}\\ 
    en Universidad de Burgos --- \@date{}\\
    Tutores: \@tutor{}\\
  }%
  \end{center}%
  \null
  \cleardoublepage
  }
\makeatother

\newcommand{\nombre}{Junhao Zhou} %%% cambio de comando

% Datos de portada
\title{Comparación de modelos de predicción de datos de temperatura y humedad procedentes de sensores IoT desplegados en un viñedo}
\author{\nombre}
\tutor{Alejandro Merino Gómez, Rubén Ruiz González}
\date{\today}

\begin{document}

\maketitle


\newpage


%%%%%%%%%%%%%%%%%%%%%%%%%%%%%%%%%%%%%%%%%%%%%%%%%%%%%%%%%%%%%%%%%%%%%%%%%%%%%%%%%%%%%%%%
\thispagestyle{empty}


\noindent\includegraphics[width=\textwidth]{cabecera}\vspace{1cm}

\noindent D. Rubén Ruiz González, profesor del departamento de Digitalización, área de Ingeniería de Sistemas y Automática.

\noindent D. Alejandro Merino Gómez, profesor del departamento de Digitalización, área de Ingeniería de Sistemas y Automática.

\noindent Exponen:

\noindent Que el alumno D. \nombre, con NIE X4146204V, ha realizado el Trabajo final de Grado en Ingeniería Informática titulado <<Comparación de modelos de predicción de datos de temperatura y humedad de suelo procedentes de sensores IoT desplegados en un viñedo>>. 

\noindent Y que dicho trabajo ha sido realizado por el alumno bajo la dirección del que suscribe, en virtud de lo cual se autoriza su presentación y defensa.

\begin{center} %\large
En Burgos, {\large \today}
\end{center}

\vfill\vfill\vfill

% Author and supervisor
\begin{minipage}{0.45\textwidth}
\begin{flushleft} %\large
Vº. Bº. del Tutor:\\[2cm]
D. Rubén Ruiz González
\end{flushleft}
\end{minipage}
\hfill
\begin{minipage}{0.45\textwidth}
\begin{flushleft} %\large
Vº. Bº. del co-tutor:\\[2cm]
D. Alejandro Merino Gómez
\end{flushleft}
\end{minipage}
\hfill

\vfill

% para casos con solo un tutor comentar lo anterior
% y descomentar lo siguiente
%Vº. Bº. del Tutor:\\[2cm]
%D. nombre tutor


\newpage\null\thispagestyle{empty}\newpage




\frontmatter

% Abstract en castellano
\renewcommand*\abstractname{Resumen}
\begin{abstract}
En el ámbito de la agricultura y el cultivo de viñedos, el análisis de datos meteorológicos desempeña un papel crucial en la toma de decisiones y la optimización de los procesos agronómicos. La recopilación y análisis de datos precisos y actualizados proporcionados por sensores permiten comprender mejor las condiciones climáticas y su impacto en las variables de interés agronómico.

\par

En este proyecto se investiga sobre distintos modelos de redes neuronales para la predicción de variables de interés agronómico en un viñedo y se realiza un análisis comparativo para detectar el mejor modelo.

\par

Se demuestra cómo el modelo que emplea una red neuronal convolucional es la que mejores resultados ha mostrado a la hora de hacer predicciones de múltiples pasos de tiempo. En cambio, dentro de los modelos probados para predicciones de un solo paso, el modelo que ha ofrecido los mejores resultados ha sido el Baseline y el LSTM.
\end{abstract}

\renewcommand*\abstractname{Descriptores}
\begin{abstract}
Minería de datos, redes neuronales, predicción, series temporales, aprendizaje automático, datos meteorológicos, preprocesamiento de datos \ldots
\end{abstract}

\clearpage

% Abstract en inglés
\renewcommand*\abstractname{Abstract}
\begin{abstract}
In the field of agriculture and vineyard cultivation, meteorological data analysis plays a fundamental role in decision making and optimization of agronomic processes. The compilation and analysis of accurate and updated data provided by sensors allows a better understanding of climatic conditions and their impact on variables of agronomic interest.

\par

This project researches different neural network models for the prediction of variables of agronomic interest in a vineyard and performs a comparative analysis to identify the best model.

\par

It is shown how the model using a convolutional neural network is the one that has shown the best results when making predictions of multiple time steps. On the other hand, within the models tested for single-step predictions, the model that has offered the best results has been the Baseline model and the LSTM model.

\end{abstract}

\renewcommand*\abstractname{Keywords}
\begin{abstract}
Data mining, neural networks, prediction, time series, machine learning, meteorological data, data preprocessing \ldots
\end{abstract}

\clearpage

% Indices
\tableofcontents

\clearpage

\listoffigures

\clearpage

\listoftables
\clearpage

\mainmatter
\capitulo{1}{Introducción}

El análisis de datos meteorológicos cumple un papel fundamental en la toma de decisiones y en la optimización de los procesos agronómicos en la agricultura y la viticultura. La recopilación y el análisis de datos precisos y actualizados que afectan a la agricultura y a la viticultura, como los datos meteorológicos, procedentes de sensores desplegados por el terreno, nos permite tener una mejor comprensión sobre las condiciones climáticas y el impacto que tiene en las variables de interés agronómico.

\par

El objetivo principal de este trabajo es realizar un análisis detallado de los datos meteorológicos recogidos por los sensores desplegados estratégicamente en un viñedo, así como la creación y comparación de diferentes modelos de renes neuronales aplicados a la predicción de series temporales relacionadas con estos datos recopilados. Específicamente, se busca utilizar estos modelos para predecir variables de interés agronómico en el viñedo lo que puede contribuir a una toma de decisiones más informada y precisa en el ámbito vitivinícola.

\par

Los sensores desplegados se encuentran en un viñedo a las afueras de Aranda de Duero, ubicado en el terreno de Burgos.

\par

Los datos recopilados a través de los sensores se han almacenado en archivos CSV independientes para cada sensor. Estos datos incluyen información relevante tomada cada cinco minutos, como el timestamp (marca de tiempo), fecha, nivel de batería, temperatura exterior, humedad exterior, temperatura a poca profundidad calibrada, humedad a poca profundidad calibrada, temperatura a mayor profundidad calibrada y humedad a mayor profundidad calibrada. Además, también se almacenan los datos de la humedad a poca profundidad y a mayor profundidad sin calibrar.

\par

Para validar los datos obtenidos por los sensores se ha llevado a cabo la obtención de datos meteorológicos históricos de la zona a través de una API y se ha realizado una comparación de los datos de cada sensor con los de los demás sensores para detectar valores anómalos.

\par

En este proyecto, se explorarán diferentes modelos para realizar predicciones de los valores meteorológicos, teniendo en cuenta la variabilidad de las condiciones climáticas. Se evaluarán y compararán distintas redes neuronales, buscando identificar aquellos que proporcionen las predicciones más precisas y confiables.
\capitulo{2}{Objetivos del proyecto}


Este proyecto está orientado al desarrollo de modelos de aprendizaje automático para las predicciones de series temporales, específicamente para la predicción de datos meteorológicos de interés agronómico en un viñedo en base a datos recogidos por sensores IoT desplegados en este.

\par

Teniendo en cuenta lo anterior, podemos establecer como objetivos de este proyecto los siguientes:

\section{Objetivos técnicos}

\begin{itemize}
    \item Realizar predicciones de variables de interés agronómico en un viñedo.
    \item Usar regresión mediante redes neuronales para la predicción en series temporales.
    \item Realizar minería de datos sobre el conjunto de datos propuesto para tratarlos previamente al desarrollo de redes neuronales.
    \item Comprender el funcionamiento de los distintos modelos de predicción a desarrollar.
    \item Desarrollar múltiples modelos de predicción empleando distintos tipos de redes neuronales.
    \item Realizar un análisis comparativo de los modelos desarrollados con los modelos de otros proyectos relacionados.
    \item Emplear un sistema de control de versiones para gestionar las modificaciones realizadas sobre el código.
    \item Realizar una investigación y una evaluación de las herramientas empleadas.
\end{itemize}

\section{Objetivos personales}

\begin{itemize}
    \item Investigar y obtener conocimientos relacionados con la minería de datos y el modelado de redes neuronales orientadas a la realización de predicciones de series temporales.
    \item Reforzar los conocimientos previos al proyecto relacionados con la minería de datos y la computación neuronal.
\end{itemize}
\setcounter{secnumdepth}{4}

\capitulo{3}{Conceptos teóricos}

\section{Minería de datos}

La minería de datos \cite{amazon:datamining} consiste en la aplicación de distintas técnicas, como inteligencia artificial o estadística, sobre grandes cantidades de datos con el objetivo de descubrir tendencias, patrones o relaciones ocultas.

\par

Si bien la minería de datos no es una técnica nueva, su relevancia ha aumentado significativamente en la época actual. Esto se debe al enorme incremento en la cantidad de datos que se generan cada día y al bajo número de analistas de datos. Debido a lo mencionado anteriormente y a que solo un pequeño porcentaje de los datos tiene un uso, la minería de datos cada vez es más importante en todos los ámbitos.

\par

La minería de datos consta de varios procesos para la realización de esta, pero una de las más conocidas es el modelo CRISP-DM~\cite{ibm:crisp-dm}. Este modelo consta de seis fases (ver Figura~\ref{fig:crisp_process}):

\begin{enumerate}
    \item Comprensión del negocio: se refiere a comprender los objetivos y los requisitos del proyecto.
    \item Compresión de los datos: esta fase está comprendida por varias tareas, que son la recopilación de los datos iniciales, la exploración de estos y la verificación de calidad de los datos.
    \item Preparación de los datos: en esta fase hay que resolver los problemas detectados en la fase anterior y mejorar el conjunto de datos.
    \item Modelado: en esta fase se seleccionan las técnicas que se van a emplear para la creación de los modelos y se genera un plan de pruebas.
    \item Evaluación: en esta fase se evalúan los modelos creados en la fase anterior y se escogen los más adecuados para el proyecto.
    \item Despliegue: en esta última fase se despliega o despliegan los modelos seleccionados para poder realizar predicciones sobre otro conjunto de datos.
\end{enumerate}

\imagen{crisp_process}{Fases del modelo CRISP-DM~\cite{imagen:crisp}}{0.7}

\section{Preprocesamiento de datos}

El preprocesamiento de datos~\cite{preprocesamiento} es una tarea fundamental en la minería de datos que nos permite obtener un formato a partir del cual podamos trabajar. Esta tarea se compone de múltiples procesos para la mejora y la limpieza del conjunto de datos. Algunos de los más importantes son el tratamiento de valores faltantes, la eliminación de datos duplicados y la eliminación de datos que se puedan considerar <<ruido>>.

\section{Normalización}

La normalización es un proceso que consiste en transformar los valores numéricos de una o más columnas para que estos valores se ajusten a una escala común con la finalidad de poder realizar comparaciones con distintos atributos con una unidad distinta y/o una escala distinta.

\par

La normalización consta de varias técnicas, las empleadas en este trabajo han sido la normalización Min-Max~\cite{normalización:minmax} y la puntuación estándar o puntuación Z~\cite{normalización:estandar}.

\par

La normalización Min-Max transforma todos los valores numéricos para que se encuentren dentro del rango [0\,,\,1]. Una desventaja de esta técnica es su sensibilidad a los valores extremos, lo que requiere un preprocesamiento de datos antes de aplicarla.

\[x_{\text{minmax}} = \frac{{x - x_{\text{min}}}}{{x_{\text{max}} - x_{\text{min}}}}\]

\par

La puntuación Z o puntuación estándar transforma los valores numéricos para que se encuentren en un rango no fijo que varía según la desviación estándar y la media. Esta técnica presenta una menor sensibilidad a los valores extremos, sin embargo, al realizar la normalización con esta técnica, cada columna tendrá un rango diferente debido a la falta de un rango fijo.

\[x_{\text{standard}} = \frac{{x - \text{mean}(x)}}{{\text{standard deviation}(x)}}\]

\section{Data windowing}

El \textit{data windowing}~\cite{datawindowing} es una técnica que consiste en dividir el conjunto de datos en subsecciones. Esta técnica se emplea principalmente para series temporales y sirve para definir dentro de cada subsección los datos que se proporcionan a la red neuronal y los datos que se quieren predecir.

\par

En el \textit{machine learning} tradicional, una mayor cantidad de datos tiende a producir mejores resultados, pero en el caso de las series temporales puede no ser así, por lo que esta técnica se suele emplear para optimizar los resultados en las predicciones de series temporales.

\imagen{windowing1}{Ejemplo de \textit{data windowing} de 6 datos proporcionados y 1 dato a predecir~\cite{imagen:data-windowing}}{0.6}

\imagen{windowing2}{Representaciones gráficas del ejemplo de \textit{data windowing}~\ref{fig:windowing1}}{0.9}

\section{Baseline}

El \textit{baseline} es un modelo de predicción que sirve como comparación con otros futuros modelos, pues no se entrena y es el más simple. Sus predicciones se basan en que la variable a predecir en un momento dado es igual que en el momento anterior. Para calcular algunas variables como la temperatura es una base razonable, ya que no varía rápidamente.

\section{Modelo lineal}

El modelo lineal es un modelo de predicción basado en la regresión lineal, un método de aprendizaje supervisado que se caracteriza por ser uno de los modelos más sencillos para realizar predicciones, ya que consta de únicamente una capa intermedia entre los datos de entrada y salida. Esta capa realiza una transformación lineal de los datos de entrada.

\section{Redes Neuronales Convolucionales}

Las Redes Neuronales Convolucionales o CNN~\cite{cnn} por sus siglas en inglés son un tipo de redes neuronales artificiales cuyo uso principal se centra en el análisis de datos con una estructura de malla, como las imágenes. Aunque su uso más común sea para reconocimiento y clasificación de imágenes y vídeos, también se pueden usar para las predicciones de series temporales.

\par

Estas redes neuronales están compuestas por tres tipos principales de capas:
\begin{itemize}
    \item Capa convolucional. 
    \item Capa de reducción de muestreo.
    \item Capa de clasificación.
\end{itemize}

La capa convolucional es la capa más importante de una red neuronal CNN, puesto que es donde ocurre la mayor parte de la computación. Esta capa se encarga de transformar los datos de entrada en una matriz bidimensional para la siguiente capa. Esa capa requiere de datos de entrada y de un filtro para poder funcionar correctamente y devuelve como salida un mapa de características. El filtro o detector de características, se moverá a través de los campos de la entrada, verificando si la característica está presente. Este proceso se conoce como convolución.

\par

El filtro es una matriz bidimensional de pesos, que representa parte de la entrada. El tamaño del filtro puede variar dependiendo del tamaño que deseemos para el campo receptivo. El filtro se aplica a un determinado área de la entrada y realiza un producto escalar entre la entrada y el filtro, este producto escalar se almacena en una matriz de salida. Tras haber realizado el producto escalar y almacenarlo, el filtro se desplaza por la entrada y repite el proceso hasta haber recorrido todos los datos de entrada. El conjunto de productos escalares almacenados en la matriz de salida se conoce como mapa de características.

\par

Tras cada operación de convolución, la red neuronal aplica una transformación de unidad lineal rectificada (ReLU) al mapa de características, lo que añade la no linealidad del modelo.

\par

La capa de reducción de muestreo tiene la función de reducir el número de parámetros de los datos de entrada para la siguiente capa. De una manera similar a la capa convolucional, la reducción de muestreo se realiza mediante la pasada de un filtro a través de los datos de entrada. Los dos principales tipos de reducción de muestreo (ver Figura~\ref{fig:pooling}) son los siguientes:
\begin{itemize}
    \item \textit{Max pooling}: Este tipo de reducción mueve el filtro a través de toda la matriz bidimensional devuelta por la capa convolucional y devuelve a la matriz de salida el valor más alto encontrado en esa instancia. Este tipo de reducción de muestreo es la más empleada de las dos.
    \item \textit{Average pooling}: Este tipo de reducción mueve el filtro a través de toda la matriz bidimensional devuelta por la capa convolucional y devuelve a la matriz de salida el valor medio de todos los valores que posea el filtro en esa instancia.
\end{itemize}

\imagen{pooling}{Ejemplo de los dos tipos de reducción de muestreo~\cite{imagen:reduccion}}{0.6}

Por último, la capa de clasificación realiza una clasificación según las características obtenidas en la fase anterior.

\par

Aunque las redes neuronales convolucionales se componen de estas capas principalmente, sería en el caso de reconocimiento y clasificación de imágenes y vídeos. En el caso de una predicción en una serie temporal, la composición de capas varía un poco, más adelante se explicará la composición de las capas empleadas para este último caso.

\section{\textit{Long short-term memory}}

Las redes \textit{Long short-term memory} (LSTM)~\cite{lstm} son un tipo de redes neuronales artificiales que están dentro del grupo de redes neuronales recurrentes. Las redes neuronales recurrentes se distinguen de las clásicas por su capacidad para <<recordar>> estados previos introduciendo bucles en el diagrama de la red.

\par

Las redes LSTM disponen de una <<memoria a largo plazo>> y de una <<memoria a corto plazo>>, a diferencia de las redes neuronales recurrentes estándar, que solo son capaces de disponer de una <<memoria a corto plazo>>. Por este motivo, son unas redes ideales para realizar predicciones de series temporales, puesto que la <<memoria a corto plazo>> ayuda a comprender y predecir los siguientes pasos de tiempo, en este caso, las siguientes horas, y la <<memoria a largo plazo>> nos ayuda a predecir las tendencias diarias o estacionales de los atributos.

\par

Una unidad LSTM está compuesta por una célula, una puerta de entrada, una puerta de salida y una puerta de <<olvido>>, como una red neuronal recurrente. La diferencia entre el modelo LSTM y un modelo RNN genérico se encuentra dentro de la célula, donde se encuentra el algoritmo de <<olvido>>.

\par

En un modelo RNN genérico, el algoritmo de conservación de información está compuesto por una sola capa sigmoidal (ver Figura~\ref{fig:RNN_Cell}), lo que determina si la información se mantiene o se <<olvida>>. 

\imagen{RNN_Cell}{Interior de una célula RNN~\cite{lstm}}{1}

En cambio, en una célula LSTM encontramos que está compuesta por cuatro capas interactuando entre ellas para determinar el mantenimiento de la información~\ref{fig:LSTM_Cell}.

\imagen{LSTM_Cell}{Interior de una célula LSTM~\cite{lstm}}{1}

Se puede encontrar más información acerca de las redes LSTM, su funcionamiento y la explicación de cada capa que se encuentra dentro de la célula en la página citada~\cite{lstm}.
\capitulo{4}{Técnicas y herramientas}

\section{Python}

Python~\cite{python} es un lenguaje de programación interpretado de alto nivel que hace énfasis en la legibilidad del código. Se emplea para el desarrollo de aplicaciones de toda clase. Posee una licencia de código abierto y es uno de los lenguajes de programación más populares. Es el lenguaje por defecto de los proyectos de \textit{machine learning}.

\par

Se ha empleado como lenguaje principal en el proyecto para el análisis de los datos y para la creación y entrenamiento de cada uno de los modelos empleados.

\section{Jupyter Notebook}

Jupyter Notebook~\cite{jupyter} es una IDE basada en la web que te permite crear documentos de Jupyter Notebook, que es un documento JSON que sigue un formato de celdas de E/S que pueden contener código.

\par

Se ha empleado para almacenar todo el código en menos ficheros y facilitar más su legibilidad.

\section{Visual Studio Code}

Visual Studio Code~\cite{vscode} es un editor de código fuente que incluye soporte para la depuración, control integrado de Git, resaltado de sintaxis y otras funciones. Es un software gratuito y de código abierto que permite también la creación y edición de documentos de Jupyter Notebook.

\par

Es el editor empleado para el desarrollo del código. Esa herramienta ha permitido la creación y la modificación de los documentos de Jupyter Notebook con facilidad, además de la ejecución de estos.

\section{MATLAB}

MATLAB~\cite{matlab} es una plataforma de programación y cálculo numérico utilizada para analizar datos, desarrollar algoritmos y crear modelos.

\par

Se ha empleado para poder obtener un resultado base para establecer como objetivo mínimo entrenando un modelo de predicciones de series temporales proporcionado por una \textit{Toolbox} de la plataforma, que es el \textit{Deep Learning Toolbox}.

\section{Bibliotecas de Python}

\subsection{Matplotlib}

Matplotlib~\cite{plt} es una biblioteca que permite la creación de gráficos bidimensionales a partir de datos contenidos en listas o \textit{arrays}.

\par

Se ha empleado para la visualización de distintos atributos de los datos empleados y para la visualización de los resultados de los distintos modelos entrenados.

\subsection{seaborn}

Seaborn~\cite{seaborn} es una biblioteca basada en \textbf{matplotlib} que permite la visualización de datos.

\par

Esta biblioteca se ha empleado específicamente para la visualización de las capas de los modelos empleados.

\subsection{Numpy}

Numpy~\cite{numpy} es una biblioteca que añade un gran número de funciones matemáticas y da soporte a la creación de matrices y vectores de gran tamaño y de múltiples dimensiones.

\par

Esta biblioteca nos ha permitido realizar distintas operaciones a lo largo del proyecto gracias a lo que ofrece.

\subsection{Pandas}

 Pandas~\cite{pandas} es una biblioteca que permite la manipulación y el análisis de datos. Ofrece herramientas para la manipulación de tablas numéricas y de series temporales.

\par

Nos ha permitido realizar el análisis y la validación del conjunto de datos empleado, así como su uso en los modelos de predicción.

\subsection{TensorFlow}

TensorFlow~\cite{tensorflow} es una biblioteca desarrollada por Google que permite la construcción y el entrenamiento de redes neuronales.

\par

Es la biblioteca empleada para el desarrollo y el entrenamiento de los modelos, debido a que es sencilla de usar en comparación con otras que también permiten la creación y el entrenamiento de modelos de aprendizaje automático, como PyTorch o Scikit-learn.

\section{APIs}

\subsection{OpenWeatherMap}

OpenWeatherMap~\cite{openweathermap} es una API que ofrece datos meteorológicos mediante una llamada a esta. También es capaz de proporcionar datos históricos y predicciones como Weatherbit. Ofrece una licencia gratuita temporal para estudiantes y docentes.

\par

Se ha empleado para la obtención de datos meteorológicos de hasta hace un año. Debido a que el conjunto de datos proporcionado dispone de un periodo superior a un año, se ha optado por buscar y emplear otra API.

\subsection{Weatherbit}

Weatherbit~\cite{weatherbit} es ora API que ofrece datos meteorológicos mediante una llamada a esta. Es capaz de proporcionar datos históricos y predicciones. Esta API también ofrece una licencia temporal gratuita para investigación.

\par

Se ha empleado como sustituto de OpenWeatherMap, pues nos ofrece la posibilidad de obtener datos meteorológicos históricos de hasta hace 20 años. Además nos ofrece la posibilidad de obtener predicciones de hasta 16 días en el futuro.

\section{Github}

GitHub~\cite{github:repo} es una plataforma en línea que permite el alojamiento de proyectos mediante el software de control de versiones Git.

Se ha empleado esta herramienta para la gestión de las versiones y de las modificaciones realizadas en el código del proyecto.

\section{\LaTeX}

\LaTeX ~\cite{wiki:latex} es un sistema de composición de textos, orientado a la creación de documentos de texto con alta calidad tipográfica. Es usado principalmente para la creación de artículos y libros científicos que incluyen, entre otras cosas, expresiones matemáticas. 

Se ha empleado para el desarrollo de la memoria y los anexos a través de una plataforma gratuita.

\section{Metodología ágil}

La metodología ágil~\cite{agile} es un conjunto de técnicas aplicadas en ciclos de trabajo cortos, con el objetivo de mejorar la eficiencia del proceso de entrega de un proyecto. Esta metodología tiene como objetivo entregar valor al cliente de manera más rápida.

\par

Scrum es un marco de la metodología ágil que se basa en la colaboración, la adaptabilidad y la entrega incremental. Se organiza en \textit{sprints}, que son periodos de tiempo predefinidos y fijos en los que se realizan las actividades de desarrollo. En este caso, se ha empleado Scrum para el desarrollo del proyecto y se ha establecido una duración de una semana o dos para los \textit{sprints}, aunque algunos han tenido periodos más largos.
\setcounter{secnumdepth}{4}

\capitulo{5}{Aspectos relevantes del proyecto}

En este apartado se va a resumir el desarrollo del proyecto y mencionar los aspectos más relevantes de este. Se va a mencionar de dónde se han obtenido los datos con los que se trabaja, cómo se han preprocesado los datos, cómo se han preparado para su uso en los modelos, la normalización de los datos, la división del conjunto de datos para poder realizar predicciones de series temporales y los modelos entrenados y la comparación realizada en base a los resultados devueltos por ellos. Se puede encontrar el desarrollo del proyecto completo en los anexos.

\section{Datos obtenidos de los sensores}

Los datos empleados en este proyecto provienen de ocho sensores IoT y de un pluviómetro desplegados en un viñedo ubicado en Burgos. Los sensores recogen datos meteorológicos del exterior y del suelo en la ubicación en el que se han desplegado. Las variables recogidas por los sensores son los siguientes:

\begin{itemize}
    \item ts (\textit{TimeStamp})
    \item fecha
    \item bateria
    \item \texttt{t\_ext} (temperatura exterior)
    \item \texttt{h\_ext} (\% humedad exterior)
    \item \texttt{t\_C\_cal} (temperatura a poca profundidad calibrada)
    \item \texttt{h\_C\_cal} (\% humedad a poca profundida calibrada)
    \item \texttt{t\_L\_cal} (temperatura a más profundidad calibrada)
    \item \texttt{h\_L\_cal} (\% humedad a más profundidad calibrada)
    \item \texttt{h\_C} (\% humedad a poca profundidad)
    \item \texttt{h\_L} (\% humedad a más profundidad)
\end{itemize}

Por último, el pluviómetro recoge el \textit{TimeStamp}, la fecha, la batería, los milímetros de precipitación acumulados desde el inicio, la variación de los milímetros de precipitación y los milímetros de precipitación por metro cuadrado.

Los datos que se han recogido a través de los sensores se han almacenado en archivos CSV independientes para cada sensor. Se han recogido datos cada cinco minutos sobre distintos parámetros relacionados con el suelo y el exterior.

\section{Obtención de datos meteorológicos a través de una API}

Este proyecto comienza con la obtención de los datos meteorológicos de la zona de los sensores para poder validar los datos obtenidos por estos últimos. Para obtener los datos, se han realizado llamadas a la API de Weatherbit introduciendo las fechas deseadas y las coordenadas. Se han hecho las llamadas necesarias para obtener datos de la localización de cada sensor y de todo el periodo en el que se han obtenido datos de los sensores.

\par

Tras haber realizado las llamadas y obtener los datos en un formato JSON, se han transformado los datos y se han almacenado en archivos CSV individuales para cada sensor.

\section{Preprocesamiento de datos}

Una vez se han obtenido los datos históricos de cada sensor, se ha realizado un preprocesamiento de los datos para comprobar la validez de cada dato que forma el conjunto de datos de cada sensor.

\section{Visualización previa en MATLAB}

Una vez concluido el preprocesamiento de los datos, se realizaron pruebas previas al desarrollo de las redes neuronales. Estas pruebas consistieron en la búsqueda de correlaciones entre las columnas que tenía el conjunto de datos y pruebas de modelos, como una regresión lineal o una predicción de series temporales.

\par

Empleando el modelo de red neuronal proporcionado por MATLAB, se ha realizado una visualización previa de los resultados deseados de la red neuronal a desarrollar. El modelo proporcionado es una red neuronal no linear autorregresiva. Para tener un objetivo base, se ha establecido que sea el valor de error devuelto por el modelo de MATLAB tras el entrenamiento.

\section{Preparación de los datos}

Antes de empezar a crear y entrenar modelos, deberemos preparar el conjunto de datos. Para ello, vamos a reducir el conjunto de datos. Puesto que el objetivo es realizar predicciones individuales para cada sensor, solo cogemos los datos de uno de los sensores, en este caso, se han escogido los datos del sensor 3, ya que al compararlo con los demás sensores, se observa que es el que tiene mayor continuidad en los datos.

\par

Tras haber escogido el conjunto de datos a emplear, vamos a reducir más aún la cantidad de datos a emplear. Para empezar, se ha transformado el conjunto de datos para que los datos no sean cada 5 minutos, sino que correspondan a un periodo de una hora, pues queremos realizar predicciones horarias y si lo hiciéramos con un conjunto de datos donde el periodo sea de 5 minutos tendríamos que realizar 12 predicciones y su precisión sería peor. Además, puesto que son datos meteorológicos, al recoger datos cada 5 minutos no se observa apenas diferencia y entraría dentro del rango de error del propio sensor.

\par

Para realizar la transformación del periodo de 5 minutos al periodo horario, se han agrupado los datos en grupos de 12 y se ha almacenado la media de cada uno de los atributos de cada uno de los grupos.

\par

Además, ya que los sensores pueden llegar a funcionar de manera incorrecta y no enviar los datos, pueden llegar a aparecer saltos y discontinuidades entre los datos. Para que eso no afecte en el entrenamiento de los modelos, se ha escogido un periodo en el que no haya ninguna discontinuidad. Ese periodo corresponde a los datos de los últimos 4 meses recogidos por el sensor.

\par

Por último, se han descartado atributos para el entrenamiento de las redes neuronales, puesto que aumentaba la complejidad y aumentaba el MSE de los modelos. Los atributos que se han empleado han sido la temperatura y la humedad bajo tierra, el mes y la hora, un total de 6 atributos.

\section{Normalización}

Para empezar a desarrollar nuestros modelos y comparar su precisión, primero debemos normalizar los valores del conjunto de datos, puesto que facilitará el aprendizaje de nuestros modelos.

\par

Para la normalización de los datos se han probado tanto la puntuación estándar como la puntuación Min-Max. A la hora de ver los resultados, se ha podido comprobar que en este caso particular es mejor la puntuación Min-Max.

\par

Una vez hemos normalizado los datos, dividimos los datos en el conjunto de datos de entrenamiento, el de validación y el de test. Para ello se ha seguido una proporción de 70/20/10 respectivamente. Para dividir el conjunto de datos se ha realizado de manera manual para que las secciones que crearemos a continuación tengan todas las horas seguidas y no haya saltos.

\section{Data windowing}

Para realizar predicciones de series temporales, debemos dividir el conjunto de datos en secciones más pequeñas y entrenar los modelos con esas secciones. Para distintos tipos de predicciones necesitaremos secciones de distintos tamaños.

\par

En este proyecto se han realizado dos tipos de predicciones, predicciones de un solo paso de tiempo y de múltiples pasos de tiempo. Para las predicciones de un solo paso se han creado secciones de 25 datos, donde se introducen como entrada los datos $t_0-t_{23}$ y se quieren predecir los datos $t_1-t_{24}$ de manera individual. En cambio, para las predicciones de múltiples pasos se han creado secciones de 48 datos, donde se introducen los datos $t_0-t_{23}$ y se quieren predecir los datos $t_{24}-t_{47}$.

\par

Para la creación de secciones de datos se ha creado una función específica para ello. Puesto que solo se realizan predicciones de uno de los parámetros del conjunto de datos, introducimos como argumento el parámetro a predecir, que en este caso es \texttt{t\_C\_cal}. Si queremos modificar la variable a predecir, debemos cambiar el argumento a la hora de crear las subsecciones.

\section{Parámetros de entrenamiento}

Para entrenar los modelos se han empleado los mismos parámetros para todos, ya que si no empleamos los mismos no obtendríamos unos resultados que podamos comparar debido a las diferencias en el entrenamiento.

\section{Modelos para predicciones de un solo paso}

Para los modelos para predicciones de un solo paso, es decir, una predicción de un solo dato en el futuro, se ha decidido que las predicciones a realizar sean, teniendo 24 horas de datos como entrada, que en este caso son 24 pasos de tiempo ($t_0-t_{23}$), que realice las predicciones de los instantes siguientes ($t_1-t_{24}$).

\imagen{single_step}{Ejemplos de secciones de datos para las predicciones de un solo paso}{1}

Para estas predicciones se han empleado tres modelos distintos, que son:

\begin{enumerate}
    \item Baseline
    \item Modelo lineal
    \item Modelo LSTM
\end{enumerate}

\subsection{Resultados}

Podemos observar en las imágenes~\ref{fig:resultados_single1} y~\ref{fig:resultados_single2} y en la tabla \ref{tabla:res-single} los resultados de los modelos de un paso.

\imagen{resultados_single1}{Resultados de los modelos Baseline y LSTM}{0.7}

\imagen{resultados_single2}{Resultados del modelo lineal}{0.7}

\tablaSmall{Resultados de los modelos de un paso}{ l l l }{res-single}{Modelo & MSE Validación & MSE Test\\}{Baseline & 0.0002 & 0.0003\\
Lineal & 0.0178 & 0.0932\\
LSTM & 0.0002 & 0.0006\\}

Podemos observar que el modelo lineal tiene un error muy alto en comparación con los otros dos modelos, teniendo en cuenta que el modelo baseline no realiza ningún cálculo y su predicción se basa en que no hay cambio y usamos este modelo como base para comparar los demás, podemos concluir que las predicciones realizadas por el modelo lineal no son acertadas y se alejan mucho del valor real. Esto se puede deber a múltiples factores, como una falta de entrenamiento, que este tipo de red neuronal no es el mejor para realizar predicciones de series temporales o que hay que modificar los parámetros del modelo para optimizarlo.

\par

En cambio, podemos ver que el modelo LSTM tiene un rendimiento muy bueno a la hora de realizar predicciones de un paso, aunque en el conjunto de datos de \textit{test} tenga un mayor error que el modelo \textit{baseline}, este error es tan pequeño que apenas es perceptible en las predicciones.

\section{Modelos para predicciones de múltiples pasos}

Las predicciones de múltiples pasos son predicciones que pueden realizarse varios pasos de tiempo en el futuro. En este caso, se han realizado predicciones de las 24 horas posteriores ($t_{24}-t_{47}$) teniendo como datos de entrada las 24 horas anteriores ($t_0-t_{23}$).

\par

Para estas predicciones se han entrenado los modelos 10 veces y se ha almacenado la media de las ejecuciones para obtener unos resultados más fiables y menos afectadas por la aleatoriedad.

\imagen{multi_step}{Ejemplos de secciones para las predicciones de múltiples pasos}{0.8}

\par

Para realizar estas predicciones se han empleado y comparado cinco modelos, que son los siguientes:

\begin{enumerate}
    \item Modelo lineal
    \item Modelo denso
    \item Modelo convolucional
    \item Modelo LSTM
    \item Modelo LSTM autorregresivo
\end{enumerate}

Para las predicciones de múltiples pasos vamos a emplear el modelo lineal como modelo base para comparar el rendimiento de los demás.

\subsection{Resultados}

Podemos observar en la siguiente tabla y en la imagen los valores correspondientes al MSE de cada modelo.

\imagen{resultados_multi}{Resultados de los modelos de múltiples pasos}{0.6}

\tablaSmall{Resultados de los modelos de múltiples pasos}{ l l l }{res-multi}{Modelo & MSE Validación & MSE Test\\}{Lineal & 0.0175 & 0.0602\\
Denso & 0.0064 & 0.0324\\
Convolucional & 0.0039 & 0.0224\\
LSTM & 0.0052 & 0.0425\\
LSTM autorregresivo & 0.0056 & 0.0419\\}

Se puede apreciar cómo los modelos más simples y el más complejo tienen un peor rendimiento respecto a los modelos que se encuentran entre ellos. Ya que los modelos lineal y denso solo realizan las predicciones en base al último paso de tiempo que se introduce como entrada, cabe esperar que sus predicciones tampoco sean tan exactas como las de los demás modelos.

\par

También podemos observar que el modelo LSTM autorregresivo tiene una precisión parecida comparada con su modelo más simple para realizar predicciones de múltiples pasos. Esto se puede deber a que un modelo demasiado complejo disponga de un número de pesos demasiado elevado para el entrenamiento y haga predicciones no tan acertadas y no sigan una tendencia regular.

\par

Cabe mencionar que todos estos modelos no son capaces de realizar una predicción con precisión cuando las temperaturas del día a predecir varían mucho respecto a las del día anterior, puesto que como solo introducimos 24 horas, las predicciones realizadas por los modelos seguirán la tendencia del día anterior y no se saldrán casi de los valores máximos o mínimos del día anterior.
\capitulo{6}{Trabajos relacionados}

Las predicciones en series temporales es un área que tiene numerosos trabajos y proyectos, puesto que se puede aplicar a toda clase de ámbitos, desde predicciones del clima hasta predicciones para criptomonedas.

\par

Debido a que este tipo de predicciones se pueden emplear en un ámbito tan grande, existe una gran variedad de trabajos a los que podemos acceder vía Internet.

\section{Kaggle}

En Kaggle podemos encontrar una variedad de trabajos relacionados con las predicciones de series temporales.

\subsection{Deep Learning for Time Series Forecasting}

En este Notebook \cite{kaggle:1} se desarrollan y se comparan distintos modelos para la predicción de las ventas diarias de distintas tiendas.

\par

En este proyecto se compara el rendimiento de 4 modelos distintos:

\begin{itemize}
    \item MLP
    \item CNN
    \item LSTM
    \item CNN-LSTM
\end{itemize}

Para los modelos, enseña todas las capas empleadas para la creación de las mismas y adjunta una explicación del modelo.

\subsection{[Tutorial] Time Series forecasting with XGBoost}

En este Notebook \cite{kaggle:2} se emplea una biblioteca de Python distinta para la predicción de series temporales, se emplea XGBoost.

\par

En este proyecto se quiere realizar una predicción del consumo de energía horario de PJM, una organización de transmisión regional de Estados Unidos.

\par

Se ha desarrollado un modelo de regresión lineal para realizar las predicciones y se ha realizado un estudio de la importancia de cada atributo que conforma el conjunto de datos para las predicciones.

\par

Adicionalmente, también se representa visualmente el día con las mejores y las peores predicciones.

\subsection{Time Series forecasting with Prophet}

En este Notebook \cite{kaggle:3} se emplea Prophet, un modelo creado por Meta, para realizar predicciones del consumo de energía horario de PJM, como en el Notebook anterior.

\par

Durante este proyecto se realiza un estudio sobre los atributos del conjunto de datos para visualizar si existe alguna tendencia.

\par

Además, para añadir un nivel de dificultad adicional a este proyecto, se realizar predicciones teniendo en cuenta el gasto de energía durante vacaciones.

\subsection{Intro to Time Series Forecasting}

Por último, otro proyecto \cite{kaggle:4} relacionado con las predicciones de series temporales es este último, que realiza una introducción al tema.

\par

En este proyecto se realizan todos los pasos anteriores a la creación y el entrenamiento de un modelo, que resulta importante para poder reducir el error y mejorar la precisión del modelo que queramos emplear.

\par

Se nos hablar de la recolección de datos, del preprocesamiento de datos, del análisis exploratorio de datos y de la validación cruzada.
\capitulo{7}{Conclusiones y Líneas de trabajo futuras}

A continuación se comentarán las conclusiones obtenidas del desarrollo de este proyecto y qué líneas de trabajo se pueden extraer de este para el futuro.

\section{Conclusiones}

Podemos concluir que, como proyecto, se han cumplido los objetivos propuestos al inicio de este. Siguiendo el modelo CRISP-DM, hemos sido capaces de desarrollar múltiples modelos de redes neuronales para la predicción de series temporales y analizar y compararlas entre sí, lo que nos ha permitido visualizar y comprender los modelos más útiles a la hora de realizar este tipo de predicciones.

\par

Además, a pesar que el conjunto de datos empleado para este proyecto no es público y es distinto de los demás proyectos que podemos encontrar en Internet, aún podemos comparar el rendimiento de los modelos entrenados por nosotros con los modelos entrenados por otros usuarios debido a que la métrica empleada para evaluar los modelos suele ser común.

\par

Además cabe recalcar, que las predicciones se han realizado sobre uno de los sensores únicamente, pero si introducimos como datos de entrada los de otros sensores, tendríamos que realizar todos los pasos del apartado 

\par

Podemos observar tanto en el modelo de ejemplo entrenado en MATLAB y en múltiples ejemplos de Internet, como los proporcionados por TensorFlow o Keras, que nuestro modelo con mejores resultados tiene una precisión bastante buena. El modelo LSTM publicado en la web de Keras como tutorial~\cite{keras:tutorial} tiene un MSE de 0.14 para el conjunto de validación. Los modelos del tutorial de TensorFlow~\cite{tensorflow:tutorial} emplean la métrica MAE, pero su equivalente en MSE es aproximadamente 0.08 para los valores con mayor precisión.

\par

Por último, también se ha cumplido con el objetivo de aprender acerca de la importancia del tratamiento de los datos previo a la creación y al entrenamiento de los modelos para reducir el error, el empleo y el funcionamiento de múltiples modelos para las predicciones de series temporales, y la importancia y el uso de las predicciones de series temporales en un ámbito real.

\section{Lineas de trabajo futuras}

Tras ver que los modelos creados y entrenados por nosotros disponen de un margen de mejora en las predicciones y que hay muchas opciones por las cuales se podría expandir este proyecto, a continuación se van a proponer unas líneas de trabajo para futuros proyectos relacionados con este:

\begin{itemize}
    \item Una expansión del análisis de modelos de predicciones de múltiples pasos con mayor complejidad. Se podría expandir la comparación de los modelos empleados desarrollando y entrenando modelos más complejos para ver su precisión y compararlos con los ya desarrollados.
    \item Modificación en las predicciones a realizar. Se podrían modificar las secciones de datos para realizar predicciones con una mayor cantidad de datos introducidos como entrada o aumentar los pasos de tiempo que se quieren predecir, por ejemplo, para realizar predicciones de 48 horas futuras.
    \item Adición de variables para realizar predicciones más precisas. Se puede realizar un proyecto donde las predicciones se realicen con una mayor cantidad de factores, como la composición del suelo donde se han desplegado los sensores, lo que nos permitiría aumentar la precisión de las predicciones.
\end{itemize}


\bibliographystyle{plain}
\bibliography{bibliografia}

\end{document}
